\documentclass{article}
\usepackage{amsmath, amssymb, amsthm}
\usepackage{hyperref}
\usepackage{geometry}
\geometry{margin=1in}

\title{Erd\H{o}s Problem \#1110}
\date{Last updated: 2025-12-07}
\author{Source: erdosproblems.com}

\begin{document}
\maketitle

\noindent\textbf{Status:} OPEN \quad \textbf{No prize}

\noindent\textbf{Tags:} number theory

\medskip

\noindent\textbf{Problem Statement:}

Let $p&#62;q\geq 2$ be two coprime integers. We call $n$ representable if it is the sum of integers of the form $p^kq^l$, none of which divide each other. If $\{p,q\}\neq \{2,3\}$ then what can be said about the density of non-representable numbers? Are there infinitely many coprime non-representable numbers?




A problem of Erd\H{o}s and Lewin \cite{ErLe96}, who proved that there are finitely many non-representable numbers if and only if $\{p,q\}=\{2,3\}$. Indeed, in \cite{Er92b} Erd\H{o}s wrote 'last year I made the following silly conjecture': every integer $n$ can be written as the sum of distinct integers of the form $2^k3^l$, none of which divide any other. He wrote 'I mistakenly thought that this was a nice and difficult conjecture but Jansen and several others found a simple proof by induction.' This simple proof is as follows: one proves the stronger fact that such a representation always exists, and moreover if $n$ is even then all the summands can be taken to be even: if $n=2m$ we are done applying the inductive hypothesis to $m$. Otherwise if $n$ is odd then let $3^k$ be the largest power of $3$ which is $\leq n$ and apply the inductive hypothesis to $n-3^k$ (which is even). Yu and Chen \cite{YuCh22} prove that the set of representable numbers has density zero whenever $q&gt;3$ or $q=3$ and $p&gt;6$ or $q=2$ and $p&gt;10$. They also prove that there are infinitely many coprime non-representable numbers if $q&gt;3$ or $q=3$ and $p\neq 5$ or $q=2$ and $p\not\in \{3,5,9\}$. Erd\H{o}s and Lewin \cite{ErLe96} also asked whether all large integers $n$ can be written as a sum of $2^k3^l$, none of which divide another, each of which is $&gt;f(n)$ for some $f(n)\to \infty$. Let $f(n)$ be the fastest growing such $f(n)$. Yu and Chen \cite{YuCh22} proved\[\frac{n}{(\log n)^{\log_23}}\ll f(n) \ll \frac{n}{\log n}.\]Yang and Zhao \cite{YaZh25} improved the lower bound to $f(n)\gg n/\log n$. van Doorn has observed in the comments that a result of Blecksmith, McCallum, and Selfridge \cite{BMS98} already implies\[f(n)\sim \frac{\log 2\log 3}{2}\frac{n}{\log n}.\]The case of three powers is the subject of [123] , and see also [845] for more on the case $\{p,q\}=\{2,3\}$. The problem [246] addresses the topic without the non-divisibility condition.




References

[BMS98] Blecksmith, Richard and McCallum, Michael and Selfridge, J. L., {$3$}-smooth representations of integers . Amer. Math. Monthly (1998), 529--543.

[Er92b] Erd\H{o}s, Paul, Some of my favourite problems in various branches of combinatorics . Matematiche (Catania) (1992), 231-240.

[ErLe96] Erd\H{o}s, P. and Lewin, Mordechai, $d$-complete sequences of integers . Math. Comp. (1996), 837-840.

[YaZh25] Yang, Quan-Hui and Zhao, Lilu, A conjecture of {Y}u and {C}hen related to the {E}rd\H os-{L}ewin theorem . Acta Arith. (2025), 277--286.

[YuCh22] Yu, Wang-Xing and Chen, Yong-Gao, On a conjecture of {E}rd\H{o}s and {L}ewin . J. Number Theory (2022), 763--778.

\noindent\textbf{Related OEIS sequences:} possible


\bigskip
\noindent\small{Source: \url{https://www.erdosproblems.com/1110}}
\end{document}
